We have, in substance, already obtained the irreducible representations of the three-dimensional rotation group. We did so while finding the eigenvalues and eigenfunctions of the total angular momentum, although without employing the ter\SourcePageJoin{470}{473}minology of group theory. Apart from a constant factor, the operators for the components of angular momentum are the operators of infinitesimal rotations\footnote{In mathematical terminology these operators are called the \emph{generators} of the rotation group.}; the angular-momentum eigenvalues therefore characterize the behaviour of wavefunctions under rotations in space. For angular momentum $j$ there are $2j+1$ distinct eigenfunctions $\psi_{jm}$, distinguished by the values $m$ of its projection and belonging to a single energy level of degeneracy $2j+1$. A rotation of the coordinate system mixes these functions among themselves, and they thus furnish an irreducible representation of the rotation group. In group-theoretic terms, therefore, the numbers $j$ label the irreducible representations of the rotation group, with one representation of dimension $2j+1$ for every $j$. Since $j$ assumes integral and half-integral values, the dimensions $2j+1$ run through all positive integers $1,2,3,\ldots$
